\documentclass[11pt,a4paper]{article}

% ---------------------------------------------------------------------------
% Relatorio Fase 1 - Correcao dos testes e Pipeline de CI
% NoFluxoUNB - FGA0314 Testes de Software
% Compilar com: pdflatex relatorio-fase1-ci.tex
% ---------------------------------------------------------------------------

\usepackage[utf8]{inputenc}
\usepackage[T1]{fontenc}
\usepackage[brazil]{babel}
\usepackage[margin=2.5cm]{geometry}
\usepackage{booktabs}
\usepackage{xcolor}
\usepackage{listings}
\usepackage{hyperref}
\usepackage{enumitem}
\usepackage{parskip}

\hypersetup{
  colorlinks=true,
  linkcolor=violet,
  urlcolor=violet,
  citecolor=violet
}

% Estilo para blocos de codigo / shell
\definecolor{codebg}{RGB}{245,245,248}
\definecolor{codecomment}{RGB}{110,110,120}
\lstset{
  backgroundcolor=\color{codebg},
  basicstyle=\ttfamily\small,
  commentstyle=\color{codecomment},
  breaklines=true,
  frame=single,
  rulecolor=\color{lightgray},
  framesep=6pt,
  xleftmargin=6pt,
  showstringspaces=false,
  columns=fullflexible,
  keepspaces=true
}

\title{\textbf{Relatório da Fase 1 --- Correção dos Testes e\\ Pipeline de Integração Contínua}}
\author{NoFluxoUNB \\ \small FGA0314 --- Testes de Software (2026.1)}
\date{\today}

\begin{document}

\maketitle

\begin{abstract}
\noindent
Este documento descreve, para a equipe, o trabalho realizado na branch
\texttt{fix/ci-pipeline}: a correção da suíte de testes Python, a reconstrução
do pipeline de Integração Contínua (CI) e a documentação associada. Explica
\emph{o que} foi alterado, \emph{por que}, e como isso concretiza a
\textbf{Fase 1} do nosso Relatório de Análise de Testes. O objetivo é que
qualquer integrante consiga entender e validar as mudanças sem precisar ler o
código linha a linha.
\end{abstract}

\tableofcontents
\bigskip

% ===========================================================================
\section{Contexto: o que era a Fase 1}
% ===========================================================================

O nosso \textbf{Relatório de Análise de Testes} diagnosticou que a base de
testes estava parcialmente quebrada e sem automação. O plano de ação definiu
três fases; a \textbf{Fase 1 (Correção Imediata)} previa três ações concretas:

\begin{enumerate}[label=\textbf{(\arabic*)}]
  \item Remapear os imports dos testes Python de \texttt{coleta\_dados} para
        \texttt{DBA/} e \textbf{remover os testes de módulos que não existem mais}.
  \item \textbf{Configurar o GitHub Actions} para rodar os testes automaticamente
        em cada Pull Request.
  \item Marcar como \emph{skip} (ou remover) testes legados irrelevantes.
\end{enumerate}

Este PR entrega exatamente esses três pontos. As Fases 2 (aumentar cobertura)
e 3 (frontend e E2E) permanecem como trabalho futuro.

% ===========================================================================
\section{Diagnóstico: por que os testes Python falhavam}
% ===========================================================================

Ao rodar a suíte, praticamente todos os arquivos Python falhavam ainda na
\textbf{importação}, com \texttt{ModuleNotFoundError}. A causa: os testes
referenciavam módulos em caminhos antigos que mudaram (ou deixaram de existir)
durante a evolução do repositório.

\begin{table}[h]
\centering
\small
\begin{tabular}{@{}lll@{}}
\toprule
\textbf{Import usado no teste} & \textbf{Situação atual} & \textbf{Ação} \\
\midrule
\texttt{coleta\_dados.scraping.*}        & migrou para \texttt{DBA/scraping}   & remapear \\
\texttt{coleta\_dados.parse\_pdf.*}      & migrou para \texttt{DBA/parse\_pdf} & remapear \\
\texttt{no\_fluxo\_backend.ai\_agent.*}  & deletado (virou backend TS)         & remover teste \\
\texttt{formatar\_para\_ragflow}         & módulo deletado                     & remover teste \\
\bottomrule
\end{tabular}
\caption{Mapeamento dos imports quebrados encontrados em \texttt{tests-python/}.}
\end{table}

A funcionalidade de IA (cliente RAGFLOW, ranking, etc.) foi reescrita em
TypeScript no \texttt{no\_fluxo\_backend} e passou a ser coberta pelo Jest.
Por isso os testes Python dessa parte ficaram \textbf{órfãos}: testavam código
que não existe mais naquele formato.

% ===========================================================================
\section{O que o fix fez}
% ===========================================================================

\subsection{Correção dos imports e robustez de \texttt{sys.path}}

Os quatro testes cujos módulos ainda existem foram remapeados:

\begin{lstlisting}[language=Python]
# Antes
from coleta_dados.scraping import extrair_turmas_sigaa as scraper
# Depois
from DBA.scraping import extrair_turmas_sigaa as scraper
\end{lstlisting}

Antes, cada arquivo inseria a raiz do projeto no \texttt{sys.path} por conta
própria (de forma frágil: um teste só funcionava porque \emph{outro} havia
ajustado o caminho antes). Centralizamos isso num \texttt{conftest.py}, que o
pytest carrega automaticamente para todos os testes:

\begin{lstlisting}[language=Python]
# tests-python/conftest.py
import os, sys
REPO_ROOT = os.path.abspath(os.path.join(os.path.dirname(__file__), ".."))
if REPO_ROOT not in sys.path:
    sys.path.insert(0, REPO_ROOT)
\end{lstlisting}

\subsection{Remoção dos testes de módulos deletados}

Foram removidos cinco arquivos cujos módulos não existem mais
(\texttt{test\_app.py}, \texttt{test\_config.py},
\texttt{test\_ragflow\_agent\_client.py},
\texttt{test\_visualizaJsonMateriasAssociadas.py} e
\texttt{test\_formatar\_para\_ragflow.py}). Manter testes que apontam para código
inexistente só gera ruído --- o histórico permanece no Git caso precisemos
recuperá-los.

\subsection{Testes obsoletos: \emph{xfail} e \emph{skip} em vez de apagar}

Depois de corrigir os imports, alguns testes passaram a \textbf{rodar}, mas a
\textbf{falhar} --- não por erro de importação, e sim porque o módulo de
produção \emph{evoluiu} e o teste ficou desatualizado. Exemplos reais:

\begin{itemize}[leftmargin=1.4em]
  \item \texttt{formatar\_turma} mudou o formato de saída (passou a usar rótulos
        \texttt{"Disciplina:"}, \texttt{"Codigo:"}, \texttt{"Ementa:"}).
  \item \texttt{salvar\_por\_departamento} ganhou um novo argumento obrigatório
        (\texttt{output\_dir}).
  \item \texttt{executar\_fluxo\_dados\_RAGFLOW} ainda usa o caminho legado
        \texttt{coleta\_dados/dados} (um bug do próprio módulo de produção).
\end{itemize}

Nesses casos, \textbf{não} reescrevemos os testes (isso é Fase 2 e exige decidir
qual é o comportamento correto). Em vez disso, marcamos cada um de forma
explícita e \textbf{rastreável}:

\begin{lstlisting}[language=Python]
@unittest.expectedFailure  # Fase 2: salvar_por_departamento agora exige output_dir
def test_salvar_por_departamento(self):
    ...
\end{lstlisting}

\paragraph{Por que \emph{xfail} e não apagar ou ignorar?}
O \texttt{@unittest.expectedFailure} (\emph{xfail}) mantém o teste sendo
\textbf{coletado e executado}; a falha é registrada como ``esperada'' e não
derruba o build. Se alguém corrigir o módulo (ou o teste) na Fase 2, o caso vira
\emph{xpass} e o CI \textbf{avisa} que aquilo já pode ser regularizado. É o
oposto de esconder o problema: é deixá-lo visível e documentado. O teste de
\texttt{upload\_pdf}, que é de \textbf{integração} (exige servidor e um PDF de
fixture), foi marcado como \texttt{skip} --- ele volta na Fase 3.

\subsection{Resultado da suíte Python}

\begin{lstlisting}[language=bash]
$ cd tests-python && python -m pytest
========== 13 passed, 1 skipped, 10 xfailed in ~3s ==========
\end{lstlisting}

A suíte agora termina com \textbf{sucesso} (código de saída 0): 13 testes
realmente passam, 1 está adiado (integração) e 10 estão marcados como dívida
conhecida da Fase 2.

% ===========================================================================
\section{O pipeline de CI reconstruído}
% ===========================================================================

\subsection{O problema anterior}

Existiam \textbf{quatro} workflows rodando testes sobrepostos
(\texttt{pipelineCI}, \texttt{all-tests}, \texttt{python-tests},
\texttt{typescript-tests}), vários deles quebrados por motivos
\emph{independentes do nosso código}:

\begin{itemize}[leftmargin=1.4em]
  \item Uso de \texttt{actions/upload-artifact@v3}, que o GitHub passou a
        \textbf{reprovar automaticamente} (descontinuado).
  \item \texttt{pytest} e \texttt{eslint} rodando na raiz do repositório, sem as
        dependências corretas (o projeto é um monorepo \texttt{pnpm}).
  \item \texttt{black} apontado como ``falho'' por diferença de versão entre
        máquinas.
\end{itemize}

\subsection{A solução: um pipeline único e legível}

Consolidamos tudo em um só arquivo,
\texttt{.github/workflows/pipelineCI.yml}, em que \textbf{cada job mapeia para um
nível da pirâmide de testes}:

\begin{table}[h]
\centering
\small
\begin{tabular}{@{}llll@{}}
\toprule
\textbf{Job} & \textbf{Papel} & \textbf{Nível} & \textbf{Ferramenta} \\
\midrule
\texttt{qualidade-python} & formatação + lint (estática) & --- & Black, Flake8 \\
\texttt{testes-python}    & backend de dados             & unidade/integração & Pytest \\
\texttt{testes-backend}   & backend TS                   & unidade & Jest \\
\texttt{testes-frontend}  & frontend Svelte              & unidade & Vitest \\
\bottomrule
\end{tabular}
\caption{Jobs do CI e seu papel na estratégia de testes.}
\end{table}

O pipeline dispara em \textbf{todo Pull Request} e em \textbf{push} para
\texttt{main}/\texttt{dev}; o PR só fica verde se todos os jobs passarem.

\subsection{Decisões importantes}

\begin{itemize}[leftmargin=1.4em]
  \item \textbf{Versões fixadas de Black e Flake8.} O CI roda em Python 3.11 e o
        Black pode formatar de modo diferente entre versões; fixar
        (\texttt{black==25.11.0}, \texttt{flake8==7.3.0}) garante que
        ``passa na minha máquina'' equivale a ``passa no CI''.
  \item \textbf{Consolidação.} Removidos \texttt{all-tests.yml},
        \texttt{python-tests.yml} e \texttt{typescript-tests.yml}; menos
        duplicação e checks mais claros.
  \item \textbf{Actions atualizadas} para \texttt{v4}/\texttt{v5} aqui e no
        \texttt{security-and-quality.yml}.
  \item \textbf{Dependências de sistema} (\texttt{tesseract-ocr},
        \texttt{poppler-utils}) instaladas no job Python por causa do
        parsing de PDF/OCR.
\end{itemize}

% ===========================================================================
\section{Branch e arquivos alterados}
% ===========================================================================

Todo o trabalho desta Fase 1 está na branch \texttt{fix/ci-pipeline}, aberta a
partir da \texttt{main}. São \textbf{19 arquivos} (excluindo este relatório):
6 adicionados, 6 modificados e 7 removidos. As tabelas abaixo indicam
\emph{onde} cada alteração foi feita.

\subsection{Arquivos adicionados}

\begin{table}[h]
\centering
\small
\begin{tabular}{@{}p{6.8cm}p{7.5cm}@{}}
\toprule
\textbf{Arquivo} & \textbf{O que é} \\
\midrule
\texttt{tests-python/conftest.py} & Coloca a raiz do repo no \texttt{sys.path} para todos os testes. \\
\texttt{docs/testes/estrategia-de-testes.md} & Visão geral da estratégia de testes (níveis, técnicas, fases). \\
\texttt{docs/testes/pipeline-ci.md} & Documentação detalhada do pipeline de CI. \\
\texttt{docs/testes/relatorio-fase1-ci.tex} & Este relatório. \\
\bottomrule
\end{tabular}
\end{table}

\subsection{Arquivos modificados}

\begin{table}[h]
\centering
\small
\begin{tabular}{@{}p{6.8cm}p{7.5cm}@{}}
\toprule
\textbf{Arquivo} & \textbf{O que mudou} \\
\midrule
\texttt{.github/workflows/pipelineCI.yml} & Reescrito: pipeline único com 4 jobs por nível; versões fixadas; actions \texttt{v4/v5}. \\
\texttt{.github/workflows/security-and-quality.yml} & \texttt{upload/download-artifact} de \texttt{v3} para \texttt{v4}. \\
\texttt{.gitignore} & Ignora artefatos de cobertura e o material de estudo local (\texttt{docs\_testes/}). \\
\texttt{mkdocs.yml} & Adiciona a seção ``Testes'' à navegação do site. \\
\texttt{tests-python/test\_converter\_json\_para\_txt.py} & Import \texttt{coleta\_dados}\,$\rightarrow$\,\texttt{DBA}; 5 testes marcados \emph{xfail}. \\
\texttt{tests-python/test\_executar\_fluxo\_dados\_RAGFLOW.py} & Corrige o alvo do \texttt{@patch}; 2 testes marcados \emph{xfail}. \\
\texttt{tests-python/test\_extrair\_turmas\_sigaa.py} & Import \texttt{coleta\_dados}\,$\rightarrow$\,\texttt{DBA}; 3 testes marcados \emph{xfail}. \\
\texttt{tests-python/test\_upload\_pdf.py} & Import remapeado; teste de integração marcado \emph{skip}. \\
\bottomrule
\end{tabular}
\end{table}

\subsection{Arquivos removidos}

\begin{table}[h]
\centering
\small
\begin{tabular}{@{}p{6.8cm}p{7.5cm}@{}}
\toprule
\textbf{Arquivo} & \textbf{Motivo} \\
\midrule
\texttt{.github/workflows/all-tests.yml} & Workflow redundante (consolidado no \texttt{pipelineCI.yml}). \\
\texttt{.github/workflows/python-tests.yml} & Workflow redundante. \\
\texttt{.github/workflows/typescript-tests.yml} & Workflow redundante. \\
\texttt{tests-python/test\_app.py} & Testa módulo deletado (\texttt{ai\_agent.app}). \\
\texttt{tests-python/test\_config.py} & Testa módulo deletado (\texttt{ai\_agent.config}). \\
\texttt{tests-python/test\_ragflow\_agent\_client.py} & Testa módulo deletado (\texttt{ai\_agent}). \\
\texttt{tests-python/test\_visualizaJsonMateriasAssociadas.py} & Testa módulo deletado (\texttt{ai\_agent}). \\
\texttt{tests-python/test\_formatar\_para\_ragflow.py} & Testa módulo deletado (\texttt{formatar\_para\_ragflow}). \\
\bottomrule
\end{tabular}
\end{table}

% ===========================================================================
\section{Como rodar localmente}
% ===========================================================================

\begin{lstlisting}[language=bash]
# Qualidade Python
pip install black==25.11.0 flake8==7.3.0
black --check .
flake8 .

# Testes Python
pip install -r tests-python/requirements.txt pytest pytest-cov pytest-mock
cd tests-python && python -m pytest

# Testes backend (TS)
cd no_fluxo_backend && npm ci && npm test

# Testes frontend (Svelte)
cd no_fluxo_frontend_svelte && npm ci && npx vitest run --passWithNoTests
\end{lstlisting}

% ===========================================================================
\section{Relação com a disciplina e próximos passos}
% ===========================================================================

A Integração Contínua \textbf{não é uma técnica de teste} (como caixa-branca ou
caixa-preta); é o \textbf{portão de qualidade} que executa, automaticamente e a
cada mudança, os testes que projetamos com essas técnicas. Ela dá sentido
prático à pirâmide de testes: garante que os testes de unidade, integração e
sistema continuem rodando e protejam o projeto de regressões.

\paragraph{Fase 2 (próxima).} Reescrever os 10 testes marcados como
\emph{xfail} conforme o comportamento atual dos módulos e \textbf{aumentar a
cobertura}, com prioridade para o \texttt{fluxograma\_controller} e o
\texttt{assistente\_controller} (backend TS, via Jest e técnicas de
caixa-branca). Há também uma limitação de cobertura a corrigir: o
\texttt{pytest.ini} mede \texttt{--cov=.} (a cobertura dos próprios arquivos de
teste), e não dos módulos de produção --- isso deve ser retargetado para
\texttt{DBA/}.

\paragraph{Fase 3.} Suíte Vitest de componentes, testes de integração com
Supabase local (Docker) e fluxo E2E com Playwright.

\bigskip
\hrule
\bigskip
\noindent\small
Branch: \texttt{fix/ci-pipeline} \quad$\cdot$\quad
Documentação complementar: \texttt{docs/testes/estrategia-de-testes.md} e
\texttt{docs/testes/pipeline-ci.md}.

\end{document}
